\documentclass[a4paper]{article}

\usepackage[T1]{fontenc}
\usepackage[utf8]{inputenc}
\usepackage{lmodern}

\usepackage[french]{babel}
\usepackage{graphicx}
\usepackage{xcolor}%				for colors
\usepackage{comment}%				for comment
\usepackage{upquote}%				for quote ''

\usepackage[toc,page]{appendix}

\usepackage{verbatim}%				Insertion de code
\usepackage{framed}%				Pour boite multipage
\usepackage{moreverb}%				Pour extension verbatim

\usepackage{listingsutf8}%				pour le code informatique
\definecolor{mGreen}{rgb}{0,0.6,0}
\definecolor{mGray}{rgb}{0.5,0.5,0.5}
\definecolor{mPurple}{rgb}{0.58,0,0.82}
\definecolor{backgroundColour}{rgb}{0.95,0.95,0.92}

\lstset{
	tabsize=4,
    extendedchars=true,
    breaklines=true,
    keepspaces=true,
	inputencoding=utf8,
    extendedchars=true,
    literate=%
                {é}{{\'e}}{1}%
                {è}{{\`e}}{1}%
                {à}{{\`a}}{1}%
                {ç}{{\c{c}}}{1}%
                {œ}{{\oe}}{1}%
                {ù}{{\`u}}{1}%
                {É}{{\'E}}{1}%
                {È}{{\`E}}{1}%
                {À}{{\`A}}{1}%
                {Ç}{{\c{C}}}{1}%
                {Œ}{{\OE}}{1}%
                {Ê}{{\^E}}{1}%
                {ê}{{\^e}}{1}%
                {î}{{\^i}}{1}%
                {ô}{{\^o}}{1}%
                {û}{{\^u}}{1}%
                {ë}{{\¨{e}}}1
                {û}{{\^{u}}}1
                {â}{{\^{a}}}1
                {Â}{{\^{A}}}1
                {Î}{{\^{I}}}1
}

\lstdefinestyle{PythonStyle}{
    backgroundcolor=\color{backgroundColour},
    commentstyle=\color{mGreen},
    keywordstyle=\color{magenta},
    numberstyle=\tiny\color{mGray},
    stringstyle=\color{mPurple},
    basicstyle=\footnotesize,
    breakatwhitespace=true,
    breaklines=true,
    captionpos=b,
    numbers=left,
    numbersep=5pt,
    language=Python
}

\usepackage{hyperref}%				for HyperText
\hypersetup{
    colorlinks=true,
    linkcolor=black,
    filecolor=magenta,
    urlcolor=cyan,
}
\urlstyle{same}

\title{Projet réseau  TM1A}
\author{AMEEUW Vincent\and CERUTTI Marc}
\date{\today}

%**************************************************************************



\begin{document}


\makeatletter
  \begin{titlepage}
  \centering
      {\large \textsc{Université de Bordeaux}}\\
      \textsc{Licence informatique}\\
    \vspace{1cm}
      \includegraphics[width=0.5\textwidth]{IMG_Latex/ubx-logo.png}\\
	\vspace{2cm}

\begin{center}
	{\large\textbf{	\@date}}\\
	\vspace{15mm}

    {\Huge   Rapport\\
    \vspace{5mm}
    \textbf{\@title}\\}

    \vspace{15mm}
    {\large AMEEUW Vincent \\ CERUTTI Marc} \\

    \vspace{15mm}
	\begin{abstract}
	\center
	Rapport pour le projet de l'enseignement \\'4TIN401U - Réseaux Info L2' (2017 - 2018) sur la mise en réseau du jeu \textit{Bomberman}\\
	\end{abstract}
\end{center}


  \end{titlepage}
\makeatother


%Table
\cleardoublepage
\tableofcontents

\newpage

%DOC
\part{Préambule}

Dans le cadre de l'enseignement '4TIN401U - Réseaux Info L2' (2017 - 2018) à l'Université de Bordeaux, en semestre 4 de Licence Informatique, nous avons dû adapter le jeu \textit{Bomberman} fait grâce à la bibliothèque Pygame en multijoueur (Description en \ref{moodle}).
\\

Le rendu final de fin d'année fut donc d'avoir un jeu \textit{Bomberman} fonctionnel en langage Python, avec un rapport fait sur notre travail avant \textbf{le vendredi 27 avril à 23h55}.
\\

Le principal objectif de cet enseignement était de nous familiariser sur la mise en réseau de projets informatiques.
Il nous a ainsi permis de mettre en pratique nos connaissances théoriques sur le réseau, la gestion des ports logiciels, des sockets, de l'envoi et de la réception de données ainsi que de leur traitement.
\\

Les contraintes techniques étaient de le faire à l'aide d'un serveur centralisé, qui ne réalise pas d'affichage graphique, mais maintient à jour l'état courant du jeu.
Seuls les clients sont en charge de l'interaction avec l'utilisateur (clavier et affichage graphique) et chaque client dispose d'une copie du modèle, qu'il doit maintenir à jour au travers des échanges réseaux avec le serveur.

	En d'autres termes :
	\begin{itemize}
\item Récupération par le client du modèle serveur à travers le réseau (map, fruits, players).
\item Gestion des connexions / déconnexions des joueurs.
\item Gestion des déplacements des joueurs.
\item Gestion des bombes.
\item Extension à de multiples joueurs.
\item Gestion des erreurs (mort violente d'un client, coupure réseau).
\item Ajout de bonus FUN dans le jeu, impliquant de faire du réseau.
	\end{itemize}


\begin{comment}
Vous devez rendre ici deux fichiers :

    un rapport PDF (de 5 à 10 pages max, police 12) décrivant votre projet, organisé comme cela :
        introduction rapide au projet
        vue d'ensemble fonctionnelle : de ce qui marche parfaitement, de ce qui ne marche pas parfaitement, des bonus et de ce qui reste à faire, ...
        architecture et implémentation : vous prendrez soin de décrire les choix techniques que vous avez effectués (TCP/UDP, select/thread, recv non bloquant, acquittement, ...), ainsi que de décrire votre protocole réseau au moyen d'un formalisme de votre choix (shéma, algorithme en pseudo-langage, automate, ...).
        bilan du projet
    une archive ZIP contenant tous les fichiers sources et ressources (images, ...) utiles au bon fonctionnement de votre projet, ainsi qu'un README détaillant comment lancer et jouer avec votre programme.

La date de rendu est fixée au \textbf{\textcolor{red}{vendredi 27 avril 23h55}}.


\end{comment}

\newpage
\part{Projet réseau}

	\section{Méthode de travail}

	Pour notre méthode de travail, on s'est d'abord mis d'accord sur les protocoles réseau à utiliser et le squelette du code sur papier, puis on a travaillé chacun de notre côté en adaptant le code de l'autre.

	Notre base de code était ainsi assez modulaire pour ne pas avoir de problèmes sur d'éventuelles modifications ou imprévus du code pour la suite.

	\section{Analyse du modèle}

	/**/

	\section{Algorithme et implémentation}
		\subsection{Protocoles}


		\subsection{Choix techniques}
		Nous avons choisi d'utiliser TCP, car ainsi nous évitions de perdre des données en transit nécessaires au bon déroulement du jeu. Cela permet d'assurer une synchronisation efficace entre les différents copies du jeu.

		En ce qui concerne la gestion des connexions, nous avons choisi d'utiliser select au lieu des threads, car du fait du partage de certaines données comme le model, la liste des sockets, etc... on ne voulait pas "bloquer" ces variables et les laisser disponibles pour modification à n'importe quel moment du jeu sans avoir de problème d'accès. Nous étions, de plus, familiers avec cet outil grâce au mini-chat réseau\footnote{Travail réalisé durant l'année en réseau.}.

		Pour l'envoi et la réception des données, nous avons choisi d'utiliser les sendall/recv de la bibliothèque socket (\ref{docpysoc}) en mode non bloquant, car de une, il ne fallait pas bloquer la boucle de jeu pour que tout s'actualise en temps réel, de deux, le sendall permet de garder un ordre d'envoi chronologique, et nous avons, avec recv, géré quand le buffer reçoit plusieurs instructions à la fois.

		Ce fut réalisé grâce aux formalismes d'envoi dans la class CommandNetwork, où l'on encode la commande de string vers byte array avec enc\textunderscore command, et on décode en altérant le modèle avec dec\textunderscore command, qui vérifie également si la commande existe. Les commandes réseau sont de cette forme :
		\textbf{CMD <arg1> <arg2>\textbackslash\textbackslash}
		\\

		Vous avez la liste des commandes en \ref{network.py}

	\section{Améliorations effectuées}
		\subsection{Collisions sur les bombes}
		L'un des principaux problèmes que nous avons rencontré en jouant est que les parties sont longues (il est difficile d'éliminer les autres).
		L'ajout de collisions avec les bombes permet de bloquer les joueurs adverses, les rendant plus simples à éliminer.
		Les parties sont de fait plus courtes mais avec plus d'action.
		\\

		Le principe de cette algorithme est de vérifier pour chaque bombe posée sa position et la position d'arrivée du joueur pour bloquer les mouvements. Nous avions réfléchi à d'autres implémentations, par exemple en rajoutant sur la map la bombe posée et en vérifiant uniquement sur la case d'arrivée du joueur s'il y a une bombe, mais cela demanderait beaucoup d'efforts pour un petit jeu et plus de mémoire "inutile" s'il n'y a pas de bombes. De plus, du fait du timer des bombes et du temps d'attente pour poser la suivante, le nombre de bombes sur le plateau est limité.

		\lstinputlisting[firstline=256,lastline=294, style=PythonStyle]{model.py}

		\subsection{Gestion des déconnexions}

		Il y a eu beaucoup de problèmes au niveau de la connection/déconnection des clients. En effet, au départ, dès que le client fermait (normalement) son jeu, le programme et la socket se fermaient en même temps, ce qui fait que le serveur ne faisait pas la différence entre une déconnection normale et une brutale.

		Nous avons donc changé le comportement de l'événement "quit", en envoyant un message au serveur pour se déconnecter, tout en laissant le jeu client tourner jusqu'à ce que ce soit le client qui gère l'erreur du fait de la déconnection par le serveur.
		\\
		L'avantage est que, pour la gestion des déconnexions brutales, si du coté serveur, un message vide ou une déconnexion impromptue du client se produit, on ferme la socket associée mais on garde l'avatar du joueur durant un temps déterminé. Ainsi, lors de ce scénario, si l'utilisateur se reconnecte avant le temps déterminé, au lieu de créer un nouveau personnage, il récupère l'ancien.

		L'inconvénient est que si le serveur ferme brutalement ou ferme la connection normalement, on peut pas différencier l'erreur.


	\section{Bilan et critique}

\newpage
\appendix
\part{Annexes}

\section{Doc python socket} \label{docpysoc}

https://docs.python.org/3/library/socket.html

\section{Moodle} \label{moodle}

https://moodle1.u-bordeaux.fr/course/view.php?id=3671

\section{Github} \label{github}

Professeur
\\
https://github.com/orel33/bomber
\\

Projet
\\
https://github.com/m21-cerutti/VM\textunderscore Reseau\textunderscore L2
\\

\newpage
\section{Code Source}

\subsection{Network.py} \label{network.py}
\lstinputlisting[style=PythonStyle]{network.py}%firstline=1,lastline=2,
\newpage

\subsection{Model.py} \label{model.py}
\lstinputlisting[style=PythonStyle]{model.py}
\newpage

\subsection{Bomber\textunderscore client.py} \label{bomber_client.py}
\lstinputlisting[style=PythonStyle]{bomber_client.py}
\newpage

\subsection{Bomber\textunderscore server.py} \label{bomber_server.py}
\lstinputlisting[style=PythonStyle]{bomber_server.py}
\newpage

\subsection{Code nécessaire à la compréhension du jeu}
\lstinputlisting[style=PythonStyle]{keyboard.py}
\newpage
\lstinputlisting[style=PythonStyle]{view.py}

\end{document}
